% -------------
\section{Release and Reproducibility}
\label{appendix:release}

This section supports the Limitations and the Ethics Statement: it states what we release, what we
withhold and why, when we queried the models, which script produces each table, and how the
evaluation set stays usable once it is public.

\paragraph{What We Release.} At publication, under the MIT License, we release:
\begin{itemize}[nosep, leftmargin=*]
\item the 150 templates and the generator, with the constants tables and the sources they record;
\item the evaluator: the final-answer check, milestone matching, the arithmetic check, the judge's
two prompts, and the statistical analysis;
\item the evaluation set: its 2,250 instances and the seed that generates them, which anyone can
check against the instance hashes we publish with the paper (\autoref{appendix:statistics});
\item every response of the evaluated models, at the providers' defaults, at matched settings, and
in the further experiments, with its scores, the judge's rulings, and the judged step flags;
\item the paraphrase pairs the domain experts keep, with their hashes;
\item the scripts behind every table and figure, listed below.
\end{itemize}

\paragraph{What We Withhold.} We withhold the domain experts' names and their individual label
files, to protect their privacy; the paper reports their labels as counts and agreement.
Until publication we also withhold the seed, so that the evaluation set cannot be regenerated before
then, and we cite rather than redistribute the textbooks the templates draw on.

\paragraph{Query Dates.} We queried every model through OpenRouter between 28 September and
6 October 2026 at the providers' default settings, and between 6 and 7 October 2026 at matched
settings (\autoref{appendix:decoding}).

\paragraph{Scripts.} Each table is written by a committed script from the released responses and
scores:
\begin{itemize}[nosep, leftmargin=*]
\item \texttt{paper\_results.py}: \autoref{tab:main_results}, the tables of
\autoref{appendix:results}, \autoref{tab:branch_domain}, \autoref{tab:paraphrase},
\autoref{tab:experiments}, \autoref{tab:errors}, and \autoref{tab:providers}, with the figures of
the results, from the outputs of the analysis scripts;
\item \texttt{paper\_setup.py}: \autoref{tab:models} and \autoref{tab:decoding}, with the setup's
numbers and the inference prompt;
\item \texttt{appendix\_statistics.py}: \autoref{tab:difficulty_stats},
\autoref{tab:template_variation}, and \autoref{tab:area};
\item \texttt{appendix\_listings.py}: the template listings of
\autoref{sec:appendix_template_examples};
\item \texttt{appendix\_certification.py}: \autoref{tab:planted_defects} and
\autoref{tab:expert_agreement};
\item \texttt{appendix\_evaluation.py}: \autoref{tab:scoring_settings} and
\autoref{tab:validation_results};
\item \texttt{worked\_example.py} and \texttt{coverage\_cases.py}: \autoref{box:worked_example}
and \autoref{box:coverage_cases}.
\end{itemize}
The first six scripts also check that the paper holds what they produce, and the last two recompute
every value they print from the responses.

\paragraph{Contamination.} Once the seed is public, the evaluation set can enter training data.
The generator draws a fresh evaluation set from any new seed by the same selection rule, so a later
evaluation can use instances that no published set holds, and we keep a held-out seed for a hosted
or later evaluation.
